\documentclass[a4paper,11pt]{ltjsarticle}
\usepackage{luatexja}
\usepackage{geometry}
\usepackage{longtable}
\usepackage{booktabs}
\usepackage{array}
\usepackage{enumitem}
\usepackage{xcolor}
\usepackage{hyperref}
\usepackage[most]{tcolorbox}
\geometry{left=22mm,right=22mm,top=22mm,bottom=24mm}
\hypersetup{colorlinks=true,linkcolor=black,urlcolor=black}
\setlength{\parindent}{0pt}
\setlength{\parskip}{0.65em}
\renewcommand{\contentsname}{目次}
\newtcolorbox{importantbox}[1]{colback=gray!8,colframe=black,title=#1,fonttitle=\bfseries,breakable}
\newtcolorbox{todobox}{colback=gray!4,colframe=black,title=TODO（図）,fonttitle=\bfseries,breakable}
\begin{document}
\begin{center}
{\LARGE\bfseries 16章 コンピューティング基礎}\par
\vspace{2mm}
{\large Computing Foundations の日本語まとめ
SWEBOK Guide v4.0a Chapter 16 を初学者向けに再構成}\par
\vspace{2mm}
{作成日：2026-05-06}\par
\end{center}

原典：Guide to the Software Engineering Body of Knowledge (SWEBOK Guide) v4.0a, Chapter 16 Computing Foundations。形式参考：01章 ソフトウェア要求の日本語まとめ。

\begin{importantbox}{この資料の主張}
コンピューティング基礎は、ソフトウェア技術者が設計・実装・性能・信頼性・安全性の判断をするための土台である。プログラムを書く能力だけでは、どのデータ構造を選ぶか、どの計算量が許容できるか、どの基本ソフトウェアやデータベースやネットワーク構成を選ぶかを説明できない。本章は、コンピュータの構成、データ構造とアルゴリズム、プログラミング言語、基本ソフトウェア、データベース、ネットワーク、人間要因、人工知能・機械学習を、ソフトウェア工学の判断材料として整理する章である。

\end{importantbox}

\begin{importantbox}{この資料の読み方}
専門用語は原則として日本語で示し、元の英語名を確認したい語は括弧内に併記する。初学者が前提知識なしに読めるように、各節では定義、実務での意味、判断の観点を分けて説明する。本資料は公式訳ではなく、SWEBOK 16章の内容を学習用に再構成した要約である。

\end{importantbox}

\tableofcontents
\newpage

\section{全体概要：この章で学ぶこと}

この章は、ソフトウェア工学を支えるコンピュータ科学の基礎を扱う。ソフトウェア技術者は、与えられたアルゴリズムを単にコードへ変換するだけではない。要求を読み、構成を設計し、適切なアルゴリズム、通信方式、性能目標、テスト方針、保守方法を選ぶ。そのためには、計算機、データ、プログラム、基本ソフトウェア、データベース、ネットワーク、人間、人工知能の基礎を横断的に理解する必要がある。

到達目標として、読者は次のことを説明できるようになることを目指す。

\begin{itemize}[leftmargin=2em]
\item システム、サブシステム、モジュール、構成要素の関係を説明できる。
\item コンピュータアーキテクチャとコンピュータ組織の違いを説明できる。
\item データ構造、アルゴリズム、計算量の関係を説明できる。
\item プログラミング言語、型、サブルーチン、オブジェクト指向、デバッグの基礎を説明できる。
\item 基本ソフトウェア、データベース、ネットワークがソフトウェア設計へ与える制約を説明できる。
\item 利用者と開発者の人間要因、人工知能・機械学習とソフトウェア工学の関係を説明できる。
\end{itemize}

\begin{todobox}
\textbf{16章の全体地図を入れる。}\par
\textbf{画像生成用プロンプト}\par
白背景、モノクロ、A4本文幅に収まる横長の学習用図解を作成する。中央に「コンピューティング基礎」を置き、周囲に「システム概念」「コンピュータ構成」「データ構造とアルゴリズム」「プログラミング」「基本ソフトウェア」「データベース」「ネットワーク」「人間要因」「人工知能・機械学習」を配置する。矢印で「ソフトウェア技術者の設計判断を支える知識」と示す。ラベルは日本語、細線、講義ノート風、余白を広くする。

\end{todobox}

\section{最初に必要な用語}

\begin{longtable}{>{\raggedright\arraybackslash}p{0.26\textwidth}>{\raggedright\arraybackslash}p{0.68\textwidth}}
\toprule
用語 & 初学者向けの説明 \\
\midrule
\endfirsthead
\toprule
用語 & 初学者向けの説明 \\
\midrule
\endhead
コンピューティング & 計算機を使って情報を表し、処理し、保存し、通信する活動全体である。 \\
システム & 目的を達成するために、ソフトウェア、ハードウェア、人、手順、データなどが組み合わさったものである。 \\
サブシステム & システムを分けた部分であり、独立した役割を持つ。さらにモジュールや部品へ分けられる。 \\
モジュール & 一つの責務を持つように切り分けたソフトウェアの単位である。変更や理解の単位になりやすい。 \\
コンピュータアーキテクチャ & コンピュータが外から見て何をするか、どのような構成要素を持つかを表す考え方である。 \\
コンピュータ組織 & 構成要素がどのようにつながり、命令をどのように実行するかという内部の実装構成である。 \\
データ構造 & データを効率よく扱うための並べ方、結び方、格納の仕方である。 \\
アルゴリズム & 問題を解くための手順であり、入力を受け取り、決められた処理をして出力を返す。 \\
計算量 & アルゴリズムが使う時間や記憶領域の増え方を表す尺度である。 \\
基本ソフトウェア & ハードウェアを管理し、アプリケーションの実行基盤を提供するソフトウェアである。通常はオペレーティングシステムを指す。 \\
プロトコル & ネットワーク上で情報をやり取りするための約束事である。 \\
人間要因 & 利用者や開発者の認知、行動、作業しやすさがシステムに与える影響である。 \\
\bottomrule
\end{longtable}

\section{導入：なぜコンピュータ科学の基礎が必要か}

プログラマは、与えられたアルゴリズムを命令列へ変換し、コンパイル、リンク、ロード、実行を通じて出力を得ることに重点を置く。一方、ソフトウェア技術者は、要求、アーキテクチャ、設計、性能、テスト、受け入れ、保守、工程を含む広い判断を行う。したがって、ソフトウェア技術者は、単にプログラムを書くだけでなく、なぜその構成が妥当なのかを説明できなければならない。

本章でいうコンピューティング基礎は、ソフトウェア技術者が高品質なソフトウェアを設計、構築、配備、保守するために必要なコンピュータ科学の基本概念である。以降では、個々の技術を丸暗記するのではなく、設計判断の材料として理解することを重視する。

\begin{importantbox}{重要}
コードは最終成果物の一部であるが、設計判断はコードより前に始まり、運用・保守まで影響する。データ構造、計算量、基本ソフトウェア、ネットワーク、データベースの知識は、後工程の手戻りを減らすための共通語である。

\end{importantbox}

\section{1 システムまたは解決策の基本概念}

解くべき問題は、機能要求、利用者との相互作用、性能要求、装置とのインタフェース、セキュリティ、脆弱性、耐久性、更新容易性などの観点から詳しく分析する必要がある。システムは、特定の機能を持つサブシステム、モジュール、構成要素が統合されたものである。

工学的に作られたシステムでは、部分同士の分け方が重要である。SWEBOK 16章は、サブシステムがモジュール性、凝集性、疎結合性を持つように設計されるべきだと整理する。

\begin{longtable}{>{\raggedright\arraybackslash}p{0.21\textwidth}>{\raggedright\arraybackslash}p{0.37\textwidth}>{\raggedright\arraybackslash}p{0.37\textwidth}}
\toprule
性質 & 意味 & 初学者向けの見方 \\
\midrule
\endfirsthead
\toprule
性質 & 意味 & 初学者向けの見方 \\
\midrule
\endhead
モジュール性 & 各サブシステムやモジュールが適切な大きさと境界を持つこと。 & 大きすぎる一枚岩ではなく、理解・変更・交換しやすい単位に分かれているかを見る。 \\
凝集性 & 一つの部分が一つの明確な仕事に集中していること。 & 一つのモジュールが、関係の薄い処理を抱え込みすぎていないかを見る。 \\
疎結合性 & 各部分ができるだけ独立して動くこと。 & ある部分を変えたとき、別の部分へ影響が広がりすぎないかを見る。 \\
\bottomrule
\end{longtable}

システムはソフトウェアだけでなく、ハードウェアも含む場合が多い。ハードウェアは、ソフトウェアの実行、入力と出力、性能目標を支える。ソフトウェア技術者は、技術、道具、データ構造、基本ソフトウェア、データベース、利用者インタフェース、プログラミング言語、アルゴリズムを、目的に合わせて選ぶ必要がある。

\begin{todobox}
\textbf{システム分解と三つの設計性質を示す図を入れる。}\par
\textbf{画像生成用プロンプト}\par
白背景、モノクロ、A4本文幅の概念図を作成する。大きな箱を「システム」とし、その中に「サブシステム」「モジュール」「構成要素」を階層的に配置する。横に三つの吹き出しで「モジュール性：適切な単位に分ける」「凝集性：一つの仕事に集中する」「疎結合性：影響範囲を小さくする」を示す。日本語ラベル、細線、講義ノート風。

\end{todobox}

\section{2 コンピュータアーキテクチャと組織}

コンピュータアーキテクチャは、コンピュータシステムの構成要素と、それが外から見てどのような役割を果たすかを表す。コンピュータ組織は、それらの構成要素が内部でどのように接続され、相互作用し、命令を実行するかを表す。前者が「何を提供するか」、後者が「どう実現するか」に近い。

システム設計では、入力・出力装置、必要なメモリ量、処理能力、補助プロセッサ、通信の帯域などを考え、要求を満たす構成を選ぶ。特に性能やリアルタイム性が重要なソフトウェアでは、ハードウェア構成の理解が設計品質に直結する。

\subsection{2.1 コンピュータアーキテクチャ}

典型的な計算システムは、記憶装置、入力・出力装置、中央処理装置（CPU）から成る。これらはバスと呼ばれる信号線で接続される。主なバスには、記憶位置や入出力装置を指定するアドレスバス、データを読み書きするデータバス、読み書きや割り込みなどの制御信号を送る制御バスがある。

\begin{longtable}{>{\raggedright\arraybackslash}p{0.22\textwidth}>{\raggedright\arraybackslash}p{0.42\textwidth}>{\raggedright\arraybackslash}p{0.30\textwidth}}
\toprule
バス & 役割 & 例として見るべきこと \\
\midrule
\endfirsthead
\toprule
バス & 役割 & 例として見るべきこと \\
\midrule
\endhead
アドレスバス & どの記憶位置または装置を対象にするかを指定する。 & 扱える記憶空間の上限に影響する。 \\
データバス & データの読み書きを運ぶ。 & 同時に運べるデータ幅や転送量に影響する。 \\
制御バス & 読み書き、割り込み、状態、リセットなどの制御を運ぶ。 & 装置同士の同期や制御に影響する。 \\
\bottomrule
\end{longtable}

\subsection{2.2 コンピュータアーキテクチャの種類}

\subsubsection{2.2.1 ノイマン型アーキテクチャ}

ノイマン型アーキテクチャは、命令とデータを同じ記憶空間に置き、算術論理装置、記憶装置、入力装置、出力装置、制御装置で構成される基本的な構成である。多くの汎用コンピュータの理解の出発点になる。

\subsubsection{2.2.2 ハーバード型アーキテクチャ}

ハーバード型アーキテクチャは、命令用記憶とデータ用記憶を分ける構成である。命令とデータを別に扱えるため、同時アクセスや保護に利点がある。修正版では、物理的には一つの記憶領域を用いながら、命令領域とデータ領域を分けて扱うこともある。

\subsubsection{2.2.3 命令セットアーキテクチャ}

命令セットアーキテクチャ（ISA）は、CPU がどのような命令、レジスタ、データ型、アドレス指定、入出力方式を持つかを定める抽象モデルである。代表的な分類に、縮小命令セットコンピュータ（RISC）と複合命令セットコンピュータ（CISC）がある。

\begin{longtable}{>{\raggedright\arraybackslash}p{0.17\textwidth}>{\raggedright\arraybackslash}p{0.45\textwidth}>{\raggedright\arraybackslash}p{0.32\textwidth}}
\toprule
分類 & 特徴 & 設計上の意味 \\
\midrule
\endfirsthead
\toprule
分類 & 特徴 & 設計上の意味 \\
\midrule
\endhead
RISC & 命令が比較的単純で、1命令が一つの仕事を行う。命令数は増えやすいが、命令実行は単純化しやすい。 & 汎用プロセッサでよく見られる。コンパイラ設計やパイプライン化を考えるときに重要である。 \\
CISC & 一つの命令が複数の処理をまとめて行うことがある。命令数は少なくできるが、命令実行は複雑になりやすい。 & 特定目的の処理や互換性を重視する場面で理解が必要になる。 \\
\bottomrule
\end{longtable}

\subsubsection{2.2.4 Flynnの分類}

Flynn の分類は、命令列とデータ列の数に着目して、並行・並列計算の構成を整理する。単一命令・単一データ（SISD）、単一命令・複数データ（SIMD）、複数命令・単一データ（MISD）、複数命令・複数データ（MIMD）が代表である。並列処理、配列処理、複数コア、分散計算を考える入口になる。

\begin{longtable}{>{\raggedright\arraybackslash}p{0.14\textwidth}>{\raggedright\arraybackslash}p{0.42\textwidth}>{\raggedright\arraybackslash}p{0.38\textwidth}}
\toprule
分類 & 読み方 & 直感的な例 \\
\midrule
\endfirsthead
\toprule
分類 & 読み方 & 直感的な例 \\
\midrule
\endhead
SISD & 一つの命令列が一つのデータ列を処理する。 & 古典的な逐次処理。 \\
SIMD & 一つの命令列が複数のデータ列を同時に処理する。 & 画像処理やベクトル演算。 \\
MISD & 複数の命令列が一つのデータ列を処理する。 & 一般的ではないが、冗長化や特殊処理の議論で登場する。 \\
MIMD & 複数の命令列が複数のデータ列を処理する。 & マルチコア、分散処理、クラスタ処理。 \\
\bottomrule
\end{longtable}

\subsubsection{2.2.5 システムアーキテクチャ}

システムアーキテクチャは、ハードウェア、ソフトウェア、モジュール、インタフェース、データ管理、通信を含む全体設計である。典型的には、統合型、分散型、プール型、コンバージド型などがある。統合型は一つの箱に機能を集める。分散型は計算と保存を別の装置に置く。プール型は共有資源を需要に応じて割り当てる。コンバージド型は分散型とプール型の利点を組み合わせ、敏捷性と拡張性を狙う。

\begin{todobox}
\textbf{代表的なコンピュータ構成の比較図を入れる。}\par
\textbf{画像生成用プロンプト}\par
白背景、モノクロ、A4本文幅の比較図を作成する。左に「ノイマン型：命令とデータが同じ記憶」、中央に「ハーバード型：命令記憶とデータ記憶を分離」、右に「Flynn分類：SISD/SIMD/MISD/MIMD」を置く。各構成に小さなブロック図を添える。ラベルは日本語、細線、講義ノート風。

\end{todobox}

\subsection{2.3 マイクロアーキテクチャまたはコンピュータ組織}

マイクロアーキテクチャは、命令セットが実際の装置内でどのように実現されるかを扱う。算術論理装置、記憶装置、入力・出力装置、制御装置の理解は、性能、割り込み、メモリ制約、装置制御を考える基礎である。

\begin{longtable}{>{\raggedright\arraybackslash}p{0.19\textwidth}>{\raggedright\arraybackslash}p{0.39\textwidth}>{\raggedright\arraybackslash}p{0.36\textwidth}}
\toprule
構成要素 & 役割 & ソフトウェア技術者が気にする点 \\
\midrule
\endfirsthead
\toprule
構成要素 & 役割 & ソフトウェア技術者が気にする点 \\
\midrule
\endhead
算術論理装置 & 算術演算と論理演算を行う。浮動小数点装置や専用処理装置を持つこともある。 & 数値計算、画像処理、暗号処理などの性能に影響する。 \\
記憶装置 & プログラムやデータを保存し、CPU がアクセスする。 & 容量、階層、キャッシュ、揮発性、アクセス速度が性能に影響する。 \\
入力・出力装置 & 外界とのデータの出入りを担う。 & 装置ドライバ、割り込み、メモリマップ入出力、DMA を理解する必要がある。 \\
制御装置 & 命令を解釈し、装置間のデータ移動と同期を制御する。 & 命令実行の順序、制御信号、装置の有効化・リセットに関係する。 \\
\bottomrule
\end{longtable}

\begin{todobox}
\textbf{CPU、メモリ、入出力、制御装置の関係図を入れる。}\par
\textbf{画像生成用プロンプト}\par
白背景、モノクロ、A4本文幅のブロック図を作成する。中央にCPUを置き、その中に「算術論理装置」「制御装置」「レジスタ」を配置する。周囲に「主記憶」「入力装置」「出力装置」「補助記憶」を置き、アドレスバス、データバス、制御バスを異なる線種で示す。日本語ラベル、講義ノート風。

\end{todobox}

\section{3 データ構造とアルゴリズム}

データ構造は、データを効果的に表し、保存し、探索し、更新するための形である。アルゴリズムは、そのデータに対して目的の処理を行う手順である。すべてのプログラムはデータを受け取り、処理し、出力を作るため、データ構造とアルゴリズムの選択は性能、保守性、拡張性、安全性に直接影響する。

\subsection{3.1 データ構造の種類}

データ型はデータの性質を表す属性である。基本データ型には文字、整数、実数、真偽値、ポインタなどがある。複合データ型は、複数の基本データ型または複合データ型を組み合わせる。抽象データ型は、内部表現よりも、利用者から見える値と操作によって定義される。

\begin{longtable}{>{\raggedright\arraybackslash}p{0.23\textwidth}>{\raggedright\arraybackslash}p{0.40\textwidth}>{\raggedright\arraybackslash}p{0.30\textwidth}}
\toprule
分類 & 例 & 使いどころ \\
\midrule
\endfirsthead
\toprule
分類 & 例 & 使いどころ \\
\midrule
\endhead
基本データ型 & 文字、整数、実数、真偽値、ポインタ。 & 値そのものを表す最小単位として使う。 \\
線形データ構造 & 配列、文字列、連結リスト、スタック、キュー、ハッシュ表。 & 順序や先頭・末尾の操作が重要な処理に使う。 \\
階層・非線形データ構造 & 木、二分木、B木、ヒープ、グラフ。 & 探索、関係表現、優先順位、依存関係の表現に使う。 \\
抽象データ型 & スタック、キュー、集合、辞書などを操作の集合として定義する。 & 実装を隠し、利用側は操作の意味だけを考えられる。 \\
\bottomrule
\end{longtable}

\begin{todobox}
\textbf{データ構造の分類ツリーを入れる。}\par
\textbf{画像生成用プロンプト}\par
白背景、モノクロ、A4本文幅の階層図を作成する。最上位に「データ構造」を置き、「基本データ型」「複合データ型」「抽象データ型」に分岐する。複合データ型から「線形：配列・リスト・スタック・キュー」「非線形：木・ヒープ・グラフ」に分岐する。日本語ラベル、細線、講義ノート風。

\end{todobox}

\subsection{3.2 データ構造への操作}

基本操作は、作成、読み取り、更新、削除である。これらは英語の頭文字から CRUD と呼ばれる。複合データでは、操作の前に目的の要素をたどる必要があるため、走査が重要になる。木構造では前順、中順、後順の走査、グラフでは深さ優先探索や幅優先探索などが使われる。

挿入や削除では、データ構造が守るべき規則を壊さないことが重要である。たとえば、二分探索木では左側が小さく右側が大きいという規則を守る必要がある。データベースでは、主キーや外部キーの整合性を壊してはならない。

\subsection{3.3 アルゴリズムとその属性}

アルゴリズムの良し悪しは、速さだけでは決まらない。正しく動くこと、異常入力に強いこと、保守しやすいこと、利用者や開発者が理解しやすいこと、拡張しやすいことも重要である。性能は重要な属性だが、性能だけを追うと可読性や保守性を損なうことがある。

\begin{longtable}{>{\raggedright\arraybackslash}p{0.23\textwidth}>{\raggedright\arraybackslash}p{0.70\textwidth}}
\toprule
属性 & 意味 \\
\midrule
\endfirsthead
\toprule
属性 & 意味 \\
\midrule
\endhead
正しさ & 想定された入力に対して要求された出力を返す。 \\
堅ろう性 & 異常入力や境界条件でも破綻しにくい。 \\
効率 & 必要な時間と記憶領域が少ない。 \\
単純性 & 仕組みを説明しやすく、誤実装しにくい。 \\
保守性 & 後から変更しやすい。 \\
拡張性 & 新しい条件や規模の増加に対応しやすい。 \\
\bottomrule
\end{longtable}

\subsection{3.4 アルゴリズムの計算量}

計算量は、入力の大きさが増えたときに、処理時間や記憶領域の使用量がどのように増えるかを表す。実行時間そのものはハードウェアや実装にも左右されるため、通常は入力サイズ n に対する増え方で比較する。

\subsection{3.5 計算量の測定}

\begin{longtable}{>{\raggedright\arraybackslash}p{0.19\textwidth}>{\raggedright\arraybackslash}p{0.38\textwidth}>{\raggedright\arraybackslash}p{0.38\textwidth}}
\toprule
記法 & 意味 & 初学者向けの理解 \\
\midrule
\endfirsthead
\toprule
記法 & 意味 & 初学者向けの理解 \\
\midrule
\endhead
大O記法 & 最悪の場合の上限を表す。 & これ以上には増えにくい、という粗い上限である。 \\
小o記法 & 厳密ではない上限を表す。 & 上限ではあるが、ぴったりではない。 \\
大Ω記法 & 最良の場合の下限を表す。 & 少なくともこれくらいは必要、という下限である。 \\
小ω記法 & 厳密ではない下限を表す。 & ゆるい意味での下限である。 \\
Θ記法 & 上限と下限が同程度であることを表す。 & 平均的または代表的な増え方の把握に使いやすい。 \\
\bottomrule
\end{longtable}

\begin{longtable}{>{\raggedright\arraybackslash}p{0.21\textwidth}>{\raggedright\arraybackslash}p{0.23\textwidth}>{\raggedright\arraybackslash}p{0.50\textwidth}}
\toprule
増え方 & 記法例 & 意味 \\
\midrule
\endfirsthead
\toprule
増え方 & 記法例 & 意味 \\
\midrule
\endhead
定数 & O(1) & 入力が増えても手順数がほぼ一定である。 \\
線形 & O(n) & 入力の大きさに比例して増える。 \\
二乗 & O(n\textasciicircum{}2) & 入力が倍になると処理が約4倍になる。 \\
三乗 & O(n\textasciicircum{}3) & 三重ループなどで現れやすい。 \\
指数 & O(2\textasciicircum{}n), O(n!) & 入力が少し増えるだけで急増する。 \\
対数 & O(log n), O(n log n) & 二分探索や効率的なソートで現れる。 \\
\bottomrule
\end{longtable}

\begin{todobox}
\textbf{計算量の増え方を比較するグラフを入れる。}\par
\textbf{画像生成用プロンプト}\par
白背景、モノクロ、A4本文幅の折れ線グラフを作成する。横軸を「入力サイズ n」、縦軸を「処理量」とし、O(1)、O(log n)、O(n)、O(n log n)、O(n\textasciicircum{}2)、O(2\textasciicircum{}n) の曲線を描く。指数の曲線は急上昇、対数は緩やかにする。日本語ラベル、講義ノート風、凡例付き。

\end{todobox}

\subsection{3.6 アルゴリズム設計}

アルゴリズムを選ぶときは、対象アプリケーションの目的と性能要求を考える。逐次処理か並列処理か、単一スレッドか複数スレッドか、暗号やグラフ処理のような専門領域かによって、適切な方法は変わる。代表的な設計法には、総当たり、再帰、分割統治、動的計画法、貪欲法、バックトラック、乱択法がある。

\begin{importantbox}{設計の観点}
良いアルゴリズム設計とは、速い手順を選ぶことだけではない。データ構造、入力の性質、実装のしやすさ、メモリ制約、失敗時の影響、保守担当者の理解しやすさを同時に考えることである。

\end{importantbox}

\subsection{3.7 ソート技法}

ソートは、データ項目を特定の順序に並べる処理である。代表的なソートには、バブルソート、クイックソート、マージソート、ヒープソート、基数ソート、バケットソート、木ソートなどがある。どのソートがよいかは、データの大きさ、すでにどの程度並んでいるか、メモリ制約、安定性の必要性によって変わる。

\subsection{3.8 探索技法}

探索は、データ集合の中から目的の項目を見つける処理である。線形探索は先頭から順番に調べる。二分探索は整列済みデータを半分ずつ絞る。木探索やハッシュ探索は、データ構造を工夫して高速化する。探索方法は、データが整列されているか、更新が多いか、メモリをどれだけ使えるかによって選ぶ。

\subsection{3.9 ハッシュ}

ハッシュは、任意の大きさのデータから固定長の値を作り、その値を使って表の位置を決める技法である。関数をハッシュ関数、得られる値をハッシュ値またはハッシュキーと呼ぶ。辞書、集合、キャッシュ、データベース索引、暗号処理で広く使われる。

ハッシュでは衝突、つまり異なる入力が同じハッシュ値になる問題が避けられない。線形探索法、二次探索法、別鎖法、開番地法などの衝突解決技法を理解する必要がある。暗号用途では、通常の高速なハッシュ関数ではなく、暗号学的性質を持つハッシュ関数を使う。

\begin{todobox}
\textbf{ソート・探索・ハッシュの使い分け図を入れる。}\par
\textbf{画像生成用プロンプト}\par
白背景、モノクロ、A4本文幅の判断フローを作成する。開始は「データから目的の結果を得たい」。分岐として「順序が必要ならソート」「一件を探すなら探索」「キーで高速アクセスするならハッシュ」を置く。各枝に例として「一覧表示」「ユーザー検索」「辞書・キャッシュ」を添える。日本語ラベル、講義ノート風。

\end{todobox}

\section{4 プログラミング基礎と言語}

コンピュータプログラムは、入力に対して命令列を実行し、目的の出力を得るための手順である。ソフトウェア技術者は、問題、実行環境、保守性、性能、開発者の熟練度を踏まえて、適切なプログラミング言語と実装方法を選ぶ必要がある。

\subsection{4.1 プログラミング言語の種類}

プログラミング言語には、マイクロプログラミング、機械語、アセンブリ言語、高水準言語がある。高水準言語は英語に近い記述でプログラムを書けるため、開発と保守を容易にする。代表的なパラダイムには、関数型、手続き型、オブジェクト指向、スクリプト、論理型がある。

\begin{longtable}{>{\raggedright\arraybackslash}p{0.23\textwidth}>{\raggedright\arraybackslash}p{0.42\textwidth}>{\raggedright\arraybackslash}p{0.28\textwidth}}
\toprule
分類 & 説明 & 例・注意点 \\
\midrule
\endfirsthead
\toprule
分類 & 説明 & 例・注意点 \\
\midrule
\endhead
マイクロプログラミング & プロセッサ内部で命令を実現する低い層の制御である。 & 通常の業務開発では直接触れないが、CPU の理解に関係する。 \\
機械語 & CPU が直接実行できる命令である。 & 人間が直接扱うには難しい。 \\
アセンブリ言語 & 機械語に対応する記号表現である。 & 装置制御や組込みで理解が必要になる。 \\
高水準言語 & 人間が読み書きしやすい抽象度の高い言語である。 & C、C++、Java、Python、JavaScript など。 \\
\bottomrule
\end{longtable}

C、C++、Java のようにコンパイラで実行形式を作る言語もあれば、JavaScript、Ruby、Python のように実行時に解釈される言語もある。実際には、仮想機械や実行時最適化を組み合わせる言語も多い。

\subsection{4.2 構文、意味論、型システム}

構文は、言語の文法である。コンパイラやインタプリタは、宣言、文、式、制御構造、関数などが文法に合っているかを調べる。意味論は、その文が何を意味するかである。同じ記号列でも、実行時の値や型によって意味が変わる場合がある。

型システムは、変数、式、関数などに型を割り当て、許される操作を制御する仕組みである。静的型付けでは、型がコンパイル時に確認される。動的型付けでは、型が実行時に決まる。型は、誤りの早期発見、実行時安全性、可読性、最適化に関係する。

\begin{longtable}{>{\raggedright\arraybackslash}p{0.19\textwidth}>{\raggedright\arraybackslash}p{0.52\textwidth}>{\raggedright\arraybackslash}p{0.23\textwidth}}
\toprule
型付け & 特徴 & 代表例 \\
\midrule
\endfirsthead
\toprule
型付け & 特徴 & 代表例 \\
\midrule
\endhead
静的型付け & 型が主にコンパイル時に確認される。早く誤りを見つけやすい。 & C、C++、Java など。 \\
動的型付け & 型が実行時に決まる。柔軟だが、実行して初めて分かる誤りもある。 & Python、Ruby、PHP、Perl など。 \\
\bottomrule
\end{longtable}

\begin{todobox}
\textbf{ソースコードが実行されるまでの流れを入れる。}\par
\textbf{画像生成用プロンプト}\par
白背景、モノクロ、A4本文幅の工程図を作成する。「ソースコード」から「字句解析」「構文解析」「中間表現」「最適化」「コード生成」「リンク」「ロード」「実行」へ矢印でつなぐ。横に「インタプリタでは一部工程が実行時に行われる」と注記する。日本語ラベル、講義ノート風。

\end{todobox}

\subsection{4.3 サブプログラムとコルーチン}

サブプログラムまたは関数は、大きなプログラムを小さな処理単位に分けるための部品である。重複を減らし、可読性と保守性を高める。値を返すものを関数、値を返さないものを手続きと呼ぶことがある。

引数の渡し方には、値渡し、参照渡し、名前渡し、結果渡し、結果値渡しなどがある。再帰では、サブプログラムが自分自身を呼び出す。再帰は木の処理などに自然だが、終了条件を明確にしなければならない。

コルーチンは、複数の入口や再開点を持つサブプログラムである。前回停止した場所を覚え、後でそこから再開する。非同期処理、イテレータ、協調的な処理の記述に役立つ。

\begin{todobox}
\textbf{サブルーチンとコルーチンの制御の違いを示す図を入れる。}\par
\textbf{画像生成用プロンプト}\par
白背景、モノクロ、A4本文幅の比較図を作成する。左に「通常のサブルーチン：呼び出し→実行→戻る」を縦の矢印で示す。右に「コルーチン：再開→一時停止→別の処理→再開」を複数の横矢印で示す。日本語ラベル、講義ノート風。

\end{todobox}

\subsection{4.4 オブジェクト指向プログラミング}

オブジェクト指向プログラミングは、データとそのデータに対する処理を一体として扱う考え方である。クラスは属性と操作を定義する雛形であり、オブジェクトはその実体である。たとえば「車両」というクラスから、乗用車、バス、トラックというオブジェクトを考えられる。

\begin{longtable}{>{\raggedright\arraybackslash}p{0.19\textwidth}>{\raggedright\arraybackslash}p{0.42\textwidth}>{\raggedright\arraybackslash}p{0.33\textwidth}}
\toprule
性質 & 意味 & 設計上の利点 \\
\midrule
\endfirsthead
\toprule
性質 & 意味 & 設計上の利点 \\
\midrule
\endhead
抽象化 & 利用者に必要な性質だけを見せ、詳細を隠す。 & 複雑さを減らし、利用側の理解を助ける。 \\
カプセル化 & データと操作をまとめ、不要なアクセスを制限する。 & データ破壊や誤用を減らす。 \\
継承 & 既存のクラスの性質を受け継ぎ、差分を追加する。 & 共通性を再利用できるが、依存関係を深くしすぎると理解しにくい。 \\
多態性 & 同じインタフェースで異なる型を扱う。 & 利用側を具体的な実装から切り離せる。 \\
\bottomrule
\end{longtable}

\begin{importantbox}{注意}
オブジェクト指向は、クラスを使えば自動的に良い設計になるという意味ではない。対象領域の概念、責務の分担、依存関係の管理を考える設計姿勢が必要である。

\end{importantbox}

\begin{todobox}
\textbf{オブジェクト指向の四要素を示す図を入れる。}\par
\textbf{画像生成用プロンプト}\par
白背景、モノクロ、A4本文幅の4象限図を作成する。象限に「抽象化」「カプセル化」「継承」「多態性」を配置し、それぞれに短い日本語説明を付ける。中央に「オブジェクト＝データ＋操作」を置く。講義ノート風、細線。

\end{todobox}

\subsection{4.5 分散プログラミングと並列プログラミング}

分散プログラミングは、ネットワークで接続された複数のコンピュータ上で、ソフトウェアの複数部分を実行し、共通の目的を達成する方法である。並列プログラミングは、一つまたは複数のコンピュータ上で処理を同時に進め、同じ目的を速く達成する方法である。

\begin{longtable}{>{\raggedright\arraybackslash}p{0.17\textwidth}>{\raggedright\arraybackslash}p{0.39\textwidth}>{\raggedright\arraybackslash}p{0.39\textwidth}}
\toprule
観点 & 分散プログラミング & 並列プログラミング \\
\midrule
\endfirsthead
\toprule
観点 & 分散プログラミング & 並列プログラミング \\
\midrule
\endhead
実行場所 & 複数のコンピュータや拠点に分かれる。 & 複数のプロセッサ、コア、スレッドで同時実行する。 \\
記憶 & 各コンピュータが独自の記憶を持つことが多い。 & 共有記憶または分散記憶を使う。 \\
通信 & ネットワーク通信が中心である。 & バス、共有記憶、プロセス間通信を使う。 \\
利点 & 可用性、拡張性、地理的分散に強い。 & 計算時間を短縮しやすい。 \\
難しさ & ネットワーク遅延、部分故障、一貫性が課題である。 & 同期、競合、デッドロック、負荷分散が課題である。 \\
\bottomrule
\end{longtable}

\subsection{4.6 デバッグ}

デバッグは、プログラムの誤りを見つけ、原因を突き止め、修正する活動である。代表的な誤りには、構文誤り、実行時誤り、論理誤りがある。構文誤りはコンパイラが見つけやすい。実行時誤りはゼロ除算、メモリ不足、不正なアドレス参照などで現れる。論理誤りは、プログラムが動くが期待と違う結果を出す誤りであり、最も見つけにくいことが多い。

\begin{longtable}{>{\raggedright\arraybackslash}p{0.19\textwidth}>{\raggedright\arraybackslash}p{0.39\textwidth}>{\raggedright\arraybackslash}p{0.36\textwidth}}
\toprule
誤り & 例 & 見つけ方 \\
\midrule
\endfirsthead
\toprule
誤り & 例 & 見つけ方 \\
\midrule
\endhead
構文誤り & 括弧の不足、予約語の誤用、型宣言の誤り。 & コンパイラや構文チェックで検出しやすい。 \\
実行時誤り & ゼロ除算、メモリオーバーフロー、不正アクセス。 & 境界値や異常値を含むテストで検出する。 \\
論理誤り & 条件式の向き違い、計算式の誤り、仕様解釈の誤り。 & デバッガ、ログ、テスト、レビューで原因を追う。 \\
\bottomrule
\end{longtable}

\subsection{4.7 標準と指針}

大規模なソフトウェアでは、多くの開発者が関わる。個人ごとの書き方がばらばらだと、統合、レビュー、保守が難しくなる。そのため、品質を重視する組織は、コーディング標準、命名規則、コメント規則、インデント、禁止すべき書き方、レビュー方法を定める。

\begin{itemize}[leftmargin=2em]
\item 対象システムに合う標準や指針を選ぶ。
\item CERT や MISRA のような公開または業界標準を必要に応じて参考にする。
\item 開発者が標準を理解し、実際に守れるよう教育する。
\item ツールとレビューで標準の遵守を確認する。
\item プロジェクトの経験から標準を定期的に見直す。
\end{itemize}

\section{5 基本ソフトウェア}

基本ソフトウェア、すなわちオペレーティングシステム（OS）は、コンピュータのハードウェアを管理し、アプリケーションに実行基盤を提供するソフトウェアである。OS の理解は、性能、並行性、入出力、ファイル、セキュリティ、ネットワークの設計判断に直結する。

OS には、バッチ処理、マルチプログラミング、時分割、単一利用者、複数利用者、単一タスク、マルチタスク、マルチスレッド、リアルタイム、ネットワーク、分散などの種類がある。小規模な組込みシステムでは、OS を使わない構成もあり得る。

\begin{longtable}{>{\raggedright\arraybackslash}p{0.28\textwidth}>{\raggedright\arraybackslash}p{0.66\textwidth}}
\toprule
OSの主要管理対象 & 説明 \\
\midrule
\endfirsthead
\toprule
OSの主要管理対象 & 説明 \\
\midrule
\endhead
プロセッサ管理 & プロセス、スレッド、スケジューリング、同期、デッドロックを扱う。 \\
メモリ管理 & 物理メモリ、仮想メモリ、ページング、置換、断片化を扱う。 \\
装置管理 & 入出力装置、装置ドライバ、割り込み、DMA、バッファリングを扱う。 \\
情報管理 & ファイル、ディレクトリ、アクセス制御、記憶装置、保護を扱う。 \\
ネットワーク管理 & 分散処理、同期、相互排除、時刻、通信、セキュリティを扱う。 \\
\bottomrule
\end{longtable}

\begin{todobox}
\textbf{OSが管理する資源の全体図を入れる。}\par
\textbf{画像生成用プロンプト}\par
白背景、モノクロ、A4本文幅の概念図を作成する。中央に「基本ソフトウェア（OS）」を置き、周囲に「CPU」「メモリ」「入出力装置」「ファイル」「ネットワーク」「アプリケーション」を配置する。OSから各資源へ管理矢印を描く。日本語ラベル、講義ノート風。

\end{todobox}

\subsection{5.1 プロセッサ管理}

プロセッサ管理では、プロセス、アドレス空間、スレッド、同期、ロック、スケジューリングを理解する必要がある。プロセスは実行中のプログラムの単位であり、スレッドはプロセス内の実行の流れである。複数の処理が同時に動くと、共有データへのアクセス順序が問題になる。

プロセス間通信には、メッセージ、パイプ、共有メモリ、セマフォなどがある。同期が不適切だと、競合、デッドロック、飢餓が発生する。CPU スケジューリングには、到着順、最短ジョブ優先、最短残り時間優先、優先度、ラウンドロビンなどがある。

\subsection{5.2 メモリ管理}

メモリ管理では、物理メモリ、仮想メモリ、二次記憶、記憶階層、リンク、割り当てを理解する。断片化には、未使用領域が分散する外部断片化と、割り当て単位の内部に無駄が残る内部断片化がある。ページング、セグメンテーション、デマンドページング、ページ置換、スラッシング、スワッピングは重要な概念である。

ページ置換戦略には、FIFO、NRU、LRU、MRU、LFU、MFU、セカンドチャンス、エージングなどがある。アプリケーションのメモリ使用特性が悪いと、OS の工夫だけでは性能を回復できないことがある。

\subsection{5.3 装置管理}

装置管理では、メモリマップ入出力と入出力マップ入出力、ブロック装置と文字装置、ポーリング、割り込み駆動、DMA、ブロッキング入出力と非ブロッキング入出力を理解する。装置ドライバは、ハードウェアとアプリケーションの間に立つソフトウェアである。

\subsection{5.4 情報管理}

情報管理では、ファイルシステム、記憶管理、ファイル属性、ディレクトリ構造、アクセス制御、保護、セキュリティを扱う。アクセス制御リスト、アクセス行列、権限、認証、ハードリンク、ソフトリンクなどは、ソフトウェアの安全性と運用性に関係する。

\subsection{5.5 ネットワーク管理}

ネットワーク管理では、分散オブジェクト、分散ファイルシステム、同期、時刻、分散相互排除、選挙アルゴリズム、マルチキャスト通信などを扱う。分散環境では、単一の正しい時刻を前提にしにくく、論理時計やベクトル時計のような概念が必要になる。

\begin{importantbox}{実務での注意}
基本ソフトウェアの制約は、アプリケーションの設計制約でもある。スレッド、ファイル、メモリ、ネットワーク、権限の扱いを知らずに実装すると、性能問題、データ破壊、セキュリティ事故、再現困難な障害につながる。

\end{importantbox}

\section{6 データベース管理}

データベースは、関連するデータ要素を整理して保存し、一つ以上のアプリケーションが高速かつ容易にアクセスできるようにした集合である。データ項目同士の関係はスキーマによって定義される。ソフトウェア技術者は、データ量、アクセス頻度、一貫性、可用性、キー、データモデル、保存モデル、利用者数を考え、適切なデータベースを選ぶ必要がある。

データベースには、関係データベース、NoSQL、列指向、オブジェクト指向、キー値、文書、階層、グラフ、時系列、ネットワーク型などがある。どれが最適かは、データの形、問い合わせ、更新頻度、一貫性要求によって変わる。

\subsection{6.1 スキーマ}

スキーマは、データ項目、テーブル、関係、索引、関数、手続き、ビュー、整合性規則を表す構造である。物理スキーマはファイルや保存構成の設計を説明し、論理スキーマはテーブルや関係の定義を説明する。

\begin{longtable}{>{\raggedright\arraybackslash}p{0.23\textwidth}>{\raggedright\arraybackslash}p{0.70\textwidth}}
\toprule
キー & 意味 \\
\midrule
\endfirsthead
\toprule
キー & 意味 \\
\midrule
\endhead
主キー & 行を一意に識別するためのキーである。 \\
代替キー & 主キー以外で一意性を持つ候補である。 \\
外部キー & 別の表の主キーなどを参照し、表同士を関連付ける。 \\
複合キー & 複数の列を組み合わせて一意性を作る。 \\
代理キー & 業務上の意味とは別に、システムが付与する識別子である。 \\
候補キー & 主キーになり得る一意な列または列集合である。 \\
\bottomrule
\end{longtable}

\begin{todobox}
\textbf{スキーマとキーの関係図を入れる。}\par
\textbf{画像生成用プロンプト}\par
白背景、モノクロ、A4本文幅のデータベース図を作成する。表「顧客」と表「注文」を置き、顧客IDを主キー、注文表の顧客IDを外部キーとして矢印でつなぐ。複合キーと代理キーの小さな例を注記する。日本語ラベル、講義ノート風。

\end{todobox}

\subsection{6.2 データモデルと保存モデル}

データモデルは、データ構造の論理面を表す。保存モデルは、データ構造の物理面を表す。データベースでは、高い一貫性と高い可用性を同時に達成することが難しい場合がある。

\begin{longtable}{>{\raggedright\arraybackslash}p{0.16\textwidth}>{\raggedright\arraybackslash}p{0.26\textwidth}>{\raggedright\arraybackslash}p{0.33\textwidth}>{\raggedright\arraybackslash}p{0.19\textwidth}}
\toprule
モデル & 正式名 & 特徴 & 向く用途 \\
\midrule
\endfirsthead
\toprule
モデル & 正式名 & 特徴 & 向く用途 \\
\midrule
\endhead
ACID & 原子性、一貫性、分離性、永続性 & トランザクションを厳密に扱い、一貫性を重視する。 & 金融、在庫、決済など。 \\
BASE & 基本的可用性、ソフト状態、結果整合性 & 柔軟で可用性を重視し、最終的な整合を受け入れる。 & 大規模分散、NoSQL、ログ、SNS系データなど。 \\
\bottomrule
\end{longtable}

\begin{longtable}{>{\raggedright\arraybackslash}p{0.28\textwidth}>{\raggedright\arraybackslash}p{0.66\textwidth}}
\toprule
保存モデル & 説明 \\
\midrule
\endfirsthead
\toprule
保存モデル & 説明 \\
\midrule
\endhead
直接接続記憶 & 記憶装置が処理するコンピュータに直接接続される。 \\
ネットワーク接続記憶 & ネットワーク上の記憶装置に複数のコンピュータがアクセスする。 \\
記憶領域ネットワーク & 複数サーバの記憶資源をネットワーク経由で効率よく提供する。 \\
\bottomrule
\end{longtable}

\subsection{6.3 データベース管理システム}

データベース管理システム（DBMS）は、データの保存、検索、保護、権限管理、同時アクセス、整合性維持を支援するソフトウェアである。代表的な構成要素には、データベースエンジン、データベースマネージャ、実行時データベースマネージャ、データベース言語、問い合わせ処理器、報告機能がある。

\begin{longtable}{>{\raggedright\arraybackslash}p{0.30\textwidth}>{\raggedright\arraybackslash}p{0.64\textwidth}}
\toprule
構成要素 & 役割 \\
\midrule
\endfirsthead
\toprule
構成要素 & 役割 \\
\midrule
\endhead
データベースエンジン & データの保存と検索の中核を担う。 \\
データベースマネージャ & 作成、削除、バックアップ、保守、更新などの機能を提供する。 \\
実行時データベースマネージャ & 認証、権限、同時アクセス、整合性を実行時に管理する。 \\
データベース言語 & 定義、操作、アクセス制御、トランザクション制御を記述する。 \\
問い合わせ処理器 & 利用者の問い合わせを解釈し、効率よく実行する。 \\
報告機能 & 条件に合うデータを抽出し、指定形式で提示する。 \\
\bottomrule
\end{longtable}

\subsection{6.4 関係データベース管理システムと正規化}

関係データベース管理システム（RDBMS）は、データを表として格納し、表同士の関係を扱う。従来の単純なファイル管理に比べ、冗長性、一貫性、同時アクセス、検索、権限管理で利点がある。

正規化は、データの冗長性や不整合を減らすために表を整理する過程である。正規化により表の数が増え、問い合わせ時間が増える場合もある。その場合は、要求に応じてあえて冗長性を戻す非正規化を行うことがある。

\begin{longtable}{>{\raggedright\arraybackslash}p{0.28\textwidth}>{\raggedright\arraybackslash}p{0.66\textwidth}}
\toprule
正規形 & 要点 \\
\midrule
\endfirsthead
\toprule
正規形 & 要点 \\
\midrule
\endhead
第1正規形 & 各セルが単一値を持ち、重複を減らす。 \\
第2正規形 & 第1正規形を満たし、部分関数従属性をなくす。 \\
第3正規形 & 第2正規形を満たし、推移的従属性をなくす。 \\
ボイス・コッド正規形 & 第3正規形を強め、決定項がスーパーキーであることを求める。 \\
第4正規形 & 多値従属性を排除する。 \\
第5正規形 & 情報を失わずにさらに分解できない状態を目指す。 \\
第6正規形・ドメインキー正規形 & 結合従属性やドメイン制約をより厳密に扱う。 \\
\bottomrule
\end{longtable}

\begin{todobox}
\textbf{正規化で表が分かれる様子を入れる。}\par
\textbf{画像生成用プロンプト}\par
白背景、モノクロ、A4本文幅のビフォー・アフター図を作成する。左に「注文表」に顧客名・住所・商品名が重複して入っている例を置く。右に「顧客表」「注文表」「商品表」へ分割した例を置き、主キーと外部キーで関係を示す。日本語ラベル、講義ノート風。

\end{todobox}

\subsection{6.5 構造化問い合わせ言語}

構造化問い合わせ言語（SQL）は、データベースの作成、更新、削除、検索に用いられる標準的な言語である。句、式、述語、問い合わせ、文などの構成要素を持ち、表の作成、ビューの作成、検索、集計、更新、削除、権限管理を行う。

SQL は標準化されているが、各データベース製品が持つ関数や細部の構文は異なる。実務では、静的SQL、埋め込みSQL、動的SQL、単純ビュー、複雑ビューを使い分ける必要がある。

\subsection{6.6 データマイニングとデータウェアハウス}

データウェアハウスは、複数のデータベースからデータを抽出し、分析しやすい共通の場所へ蓄積する仕組みである。データマイニングは、その蓄積データからパターンや知見を抽出する。関連分析、クラスタリング、分類、系列パターン、予測などが代表的な技法である。

\subsection{6.7 データベースのバックアップと回復}

データベースは故障や破損の影響を受ける。トランザクション更新では、コミット、取り消し、遅延更新、即時更新、キャッシュ、シャドウページングなどを考慮する。バックアップには、全体バックアップ、差分バックアップ、トランザクションログバックアップがある。

\begin{importantbox}{実務での注意}
データベース設計では、正規化だけを目的にしない。一貫性、性能、更新頻度、可用性、監査、障害回復、利用者権限を含めて判断する必要がある。

\end{importantbox}

\section{7 コンピュータネットワークと通信}

コンピュータネットワークは、情報を共有するために接続された装置の集まりである。通信、装置共有、データ保存、セキュリティ、遠隔監視、業務連携、ウェブ閲覧など、多くのソフトウェアはネットワークを前提にしている。分散計算、グリッド計算、クラウド計算もネットワーク原理に基づく。

\subsection{7.1 コンピュータネットワークの種類}

\begin{longtable}{>{\raggedright\arraybackslash}p{0.32\textwidth}>{\raggedright\arraybackslash}p{0.62\textwidth}}
\toprule
種類 & 説明 \\
\midrule
\endfirsthead
\toprule
種類 & 説明 \\
\midrule
\endhead
個人領域ネットワーク・家庭内ネットワーク & 個人または家庭内の短距離接続である。 \\
局所領域ネットワーク & 建物や組織内のネットワークである。 \\
無線局所領域ネットワーク & 無線で構成される局所領域ネットワークである。 \\
広域ネットワーク & 都市、国、拠点間など広範囲をつなぐ。 \\
キャンパス領域ネットワーク & 大学や企業キャンパス内の複数建物をつなぐ。 \\
都市領域ネットワーク & 都市規模のネットワークである。 \\
記憶領域ネットワーク & 記憶装置を高速に共有するためのネットワークである。 \\
仮想私設ネットワーク & 公衆網上に仮想的な私設通信路を作る。 \\
\bottomrule
\end{longtable}

\subsection{7.2 ネットワークの層状アーキテクチャ}

通信システムは複雑であるため、層に分けて考える。各層は特定の役割を持ち、下位層のサービスを使って上位層へ機能を提供する。同じ層同士の会話の約束をプロトコルという。層の基本要素は、サービス、プロトコル、インタフェースである。

\begin{longtable}{>{\raggedright\arraybackslash}p{0.23\textwidth}>{\raggedright\arraybackslash}p{0.70\textwidth}}
\toprule
要素 & 意味 \\
\midrule
\endfirsthead
\toprule
要素 & 意味 \\
\midrule
\endhead
サービス & ある層が隣接する上位層に提供する機能である。 \\
プロトコル & 同じ層の相手と情報を交換するための規則である。 \\
インタフェース & 隣接する層の間で情報を渡すための接点である。 \\
\bottomrule
\end{longtable}

\subsection{7.3 開放型システム間相互接続モデル}

開放型システム間相互接続モデル（OSI参照モデル）は、ネットワーク通信を七つの層に分ける参照モデルである。各層は独立した役割を持ち、上位または下位の層から受け取ったデータを処理して渡す。

\begin{longtable}{>{\raggedright\arraybackslash}p{0.14\textwidth}>{\raggedright\arraybackslash}p{0.28\textwidth}>{\raggedright\arraybackslash}p{0.52\textwidth}}
\toprule
層 & 名称 & 役割の直感 \\
\midrule
\endfirsthead
\toprule
層 & 名称 & 役割の直感 \\
\midrule
\endhead
第1層 & 物理層 & 信号、ケーブル、電波など実際の媒体を扱う。 \\
第2層 & データリンク層 & 同じリンク上でフレームを送る。 \\
第3層 & ネットワーク層 & ネットワーク間でパケットを配送する。 \\
第4層 & トランスポート層 & 端末間の通信品質や接続を扱う。 \\
第5層 & セッション層 & 通信の開始、維持、終了を扱う。 \\
第6層 & プレゼンテーション層 & 表現形式、変換、暗号化などを扱う。 \\
第7層 & アプリケーション層 & 利用者アプリケーションに近い通信機能を扱う。 \\
\bottomrule
\end{longtable}

\begin{todobox}
\textbf{OSI参照モデルと層間通信を示す図を入れる。}\par
\textbf{画像生成用プロンプト}\par
白背景、モノクロ、A4本文幅の縦長図を作成する。左右に二つのコンピュータを置き、それぞれ7層の積み重ねとして「物理」「データリンク」「ネットワーク」「トランスポート」「セッション」「プレゼンテーション」「アプリケーション」を描く。同じ層同士に点線で「プロトコル」、隣接層に「インタフェース」を示す。日本語ラベル、講義ノート風。

\end{todobox}

\subsection{7.4 カプセル化とカプセル解除}

送信側では、上位層から受け取ったデータに各層がヘッダや末尾情報を付け加える。この処理をカプセル化という。受信側では、各層が自分に関係する情報を取り出し、残りを上位層へ渡す。この処理をカプセル解除という。各層のデータ単位はプロトコルデータ単位（PDU）と呼ばれる。

\begin{todobox}
\textbf{カプセル化とカプセル解除の図を入れる。}\par
\textbf{画像生成用プロンプト}\par
白背景、モノクロ、A4本文幅の左右対称図を作成する。左側を送信側として、上位データに各層のヘッダが順に付く様子を下向きに描く。右側を受信側として、各層でヘッダが外れる様子を上向きに描く。中央に「PDU」を置く。日本語ラベル、講義ノート風。

\end{todobox}

\subsection{7.5 アプリケーション層プロトコル}

アプリケーション層は、利用者のアプリケーションにサービスとインタフェースを提供する。ファイル転送、遠隔ログイン、電子メール、名前解決、ネットワーク管理、印刷などに関するプロトコルがある。FTP、TFTP、NFS、Telnet、SMTP、DNS、SNMP、DHCP などは代表例である。

\subsection{7.6 信頼でき効率的なネットワークの設計技法}

現代の業務システムでは、ネットワークは常時稼働、信頼性、効率、拡張性を求められる。設計では、業務目標と技術解決策、拡張のロードマップを明確にする。基本目標は、信頼性、セキュリティ、可用性、管理容易性である。

設計では、ファイアウォール、LAN、VLAN、サブネット、サービス品質、非武装地帯、スパニングツリー、ポートチャネル、無線アクセスポイント、物理セキュリティ、監視、障害検知を考える。攻撃や脆弱性は常に変化するため、設計後も継続的な監視と更新が必要である。

\subsection{7.7 インターネットプロトコル群}

インターネットプロトコル群、または TCP/IP は、インターネット上で接続されたコンピュータ同士がデータを通信するための一連の規約である。OSIの上位三層はTCP/IPではアプリケーション層へまとめて扱われることが多く、IP がネットワーク層の中心的役割を担う。

代表的なプロトコルには、TCP/IP、UDP/IP、SMTP、PPP、FTP、SFTP、HTTP、HTTPS、Telnet、POP3、VoIP などがある。IPv4 は 32 ビットのアドレス、IPv6 は 128 ビットのアドレスを使う。IPv4 と IPv6 の移行では、二重スタック、トンネリング、NAT64 などを理解する必要がある。

\subsection{7.8 無線ネットワークと移動体ネットワーク}

無線ネットワークは、ケーブルなしで装置を接続し通信できるようにする。無線個人領域ネットワーク、無線局所領域ネットワーク、無線広域ネットワークなどがある。移動体ネットワークでは、セルと呼ばれる地域を基地局が担当し、隣接セルでは干渉を避けるため異なる周波数を使う。

データ伝送には、周波数分割多元接続、時分割多元接続、符号分割多元接続、空間分割多元接続などがある。1Gから5Gまでの世代差では、中核網、アクセス方式、周波数、帯域幅、通信技術が変化する。

\subsection{7.9 セキュリティと脆弱性}

無線や移動体通信は便利だが、適切に保護しなければ攻撃を受けやすい。リスクには、便乗接続、無線探索、偽アクセスポイント、盗聴、未承認アクセス、端末盗難、フィッシング系攻撃、暗号方式への攻撃、電波妨害などがある。

\begin{itemize}[leftmargin=2em]
\item 初期パスワードを変更し、定期的に更新する。
\item 許可された利用者だけにアクセスを制限する。
\item システム内とネットワーク上のデータを暗号化する。
\item 複数段のファイアウォールを利用する。
\item SSIDの公開や共有設定を慎重に扱う。
\item ウイルス対策、VPN、装置のセキュリティパッチを継続的に更新する。
\end{itemize}

\begin{importantbox}{実務での注意}
ネットワーク設計は、つながることだけを目的にしない。可用性、性能、管理、認証、暗号化、監視、障害時の切り替え、ログ取得まで含めて設計する必要がある。

\end{importantbox}

\section{8 利用者と開発者の人間要因}

ソフトウェア利用者と開発者は、異なる目的、知識、行動様式を持つ。人間とコンピュータの相互作用（HCI）は、利用者が計算システムと自然にやり取りできるようにする技術と設計を扱う。利用者満足度は、利用者体験（UX）として評価されることが多い。

\subsection{8.1 利用者の人間要因}

利用者は、堅ろうで、安全で、直感的な画面を持ち、少ない手順で目的を達成でき、応答が速く一貫しているソフトウェアを期待する。メッセージは、成功時も失敗時も明確で完全である必要がある。可能な場合、処理中断や取り消しができることも重要である。

利用者インタフェース開発では、利用者のプロフィール、利用する入出力装置、入力方法、出力方法、障害時の許容度、性能期待を確認する。多くの場合、試作から始め、利用者の反応を取り入れて反復的に改善する。

\subsection{8.2 開発者の人間要因}

ソフトウェアは、開発期間より長く使われることが多い。保守する技術者が、最初に書いた技術者と同じとは限らない。したがって、コードは他者が読むことを前提に書く必要がある。意味のある文書化、適切なコーディング標準、読みやすい命名、構造化、整ったコメントは、保守性を高める。

良いコードは、書く人の都合だけでなく、後から読む人の理解を重視する。読みやすさ、単純さ、一貫したスタイル、モジュール性は、長期保守において性能と同じくらい重要な設計品質である。

\begin{todobox}
\textbf{利用者要因と開発者要因を対比する図を入れる。}\par
\textbf{画像生成用プロンプト}\par
白背景、モノクロ、A4本文幅の2列図を作成する。左列を「利用者要因」とし、「直感的操作」「明確なメッセージ」「高速応答」「取り消し」「安全性」を置く。右列を「開発者要因」とし、「可読性」「文書化」「命名規則」「モジュール化」「保守性」を置く。中央に「人間の限界を前提に設計する」と書く。日本語ラベル、講義ノート風。

\end{todobox}

\section{9 人工知能と機械学習}

知能とは、情報や知識を獲得し、関連付け、特定の課題に対して適切な判断を行う能力である。人工知能（AI）は、計算機システムに人間のような知的振る舞いを実現させることを目指す。機械学習（ML）は、経験やデータから学習し、得た知識を使って判断する仕組みである。深層学習は、人工ニューラルネットワークを用いて学習と予測を行う方法である。

\subsection{9.1 推論}

推論は、利用可能な情報を分析し、状況の原因や結論を導く能力である。AIにおいて、推論の結果は次に何をすべきかを決める材料になる。

\begin{longtable}{>{\raggedright\arraybackslash}p{0.23\textwidth}>{\raggedright\arraybackslash}p{0.70\textwidth}}
\toprule
推論 & 説明 \\
\midrule
\endfirsthead
\toprule
推論 & 説明 \\
\midrule
\endhead
演繹推論 & 前提が真で推論規則が正しければ、結論も真になる推論である。 \\
帰納推論 & 観察事例から一般化する推論である。結論は確からしいが絶対ではない。 \\
仮説推論 & 不完全な情報から、最もありそうな説明を導く推論である。 \\
常識推論 & 過去の経験や一般常識から状況を判断する推論である。 \\
単調推論 & 一度得た結論が、新しい情報で取り消されない推論である。 \\
非単調推論 & 新しい知識により結論が変化する推論である。 \\
\bottomrule
\end{longtable}

\subsection{9.2 学習}

学習は、観察、実験、経験から得た情報を将来の判断に活用できるようにすることである。AIシステムでは、入力、出力、環境からのフィードバックを使って、予測や判断の能力を改善する。

\begin{longtable}{>{\raggedright\arraybackslash}p{0.23\textwidth}>{\raggedright\arraybackslash}p{0.50\textwidth}>{\raggedright\arraybackslash}p{0.22\textwidth}}
\toprule
学習 & 説明 & 典型例 \\
\midrule
\endfirsthead
\toprule
学習 & 説明 & 典型例 \\
\midrule
\endhead
教師あり学習 & ラベル付きデータで訓練する。 & 分類、回帰。 \\
教師なし学習 & ラベルなしデータから構造やパターンを見つける。 & クラスタリング、次元削減。 \\
半教師あり学習 & ラベル付きとラベルなしの両方を使う。 & ラベル付けコストが高い場合。 \\
強化学習 & 環境との相互作用から報酬や失敗を受け取り、試行錯誤で学ぶ。 & ゲーム、ロボット制御。 \\
\bottomrule
\end{longtable}

\subsection{9.3 モデル}

AIモデルは、関連データに基づいて判断を行う推論器またはアルゴリズムである。線形回帰、ロジスティック回帰、人工ニューラルネットワーク、決定木、ナイーブベイズ、サポートベクトル機械、ランダムフォレストなどがある。モデルごとに向くデータ、問題、説明可能性、必要データ量、計算量が異なる。

\begin{longtable}{>{\raggedright\arraybackslash}p{0.30\textwidth}>{\raggedright\arraybackslash}p{0.64\textwidth}}
\toprule
モデル & 主な用途の直感 \\
\midrule
\endfirsthead
\toprule
モデル & 主な用途の直感 \\
\midrule
\endhead
線形回帰 & 連続値を予測する。 \\
ロジスティック回帰 & 分類や確率的判断に使う。 \\
人工ニューラルネットワーク & 複雑な特徴を学習し、画像・音声・言語で使われる。 \\
決定木 & 条件分岐で判断過程を表す。 \\
ナイーブベイズ & 特徴の独立性を仮定し、分類に使う。 \\
サポートベクトル機械 & 限られたデータで分類境界を作る。 \\
ランダムフォレスト & 複数の決定木を組み合わせ、分類や回帰を行う。 \\
\bottomrule
\end{longtable}

\begin{todobox}
\textbf{学習方法と代表モデルの対応図を入れる。}\par
\textbf{画像生成用プロンプト}\par
白背景、モノクロ、A4本文幅のマップ図を作成する。左に「教師あり」「教師なし」「半教師あり」「強化学習」を置き、右に「回帰」「分類」「クラスタリング」「意思決定」「制御」を置く。代表モデルとして線形回帰、決定木、ニューラルネットワーク、ランダムフォレストを小さなタグで配置する。日本語ラベル、講義ノート風。

\end{todobox}

\subsection{9.4 知覚と問題解決}

AIにおける問題解決は、利用者の命令や環境から得た情報を理解し、知識ベースや推論機構を使って適切な行動を選ぶことである。外界を扱うAIシステムは、カメラ、マイク、温度、圧力、光などのセンサからデータを得て、分析し、行動する。

AIシステムは、特定作業を知的に行う型、現在の状況に反応する型、自己認識や心の理論を想定する型などに整理されることがある。実務では、こうした分類を万能な成熟度表としてではなく、AIがどの範囲の状況を扱えるかを考える補助線として使うべきである。

\subsection{9.5 自然言語処理}

自然言語処理（NLP）は、人間が日常的に使う言語をAIシステムが理解し、命令や質問に応答できるようにする分野である。音声命令を扱う場合は、単語だけでなく、発音、言い回し、俗語、文脈を理解する必要がある。

\subsection{9.6 AIとソフトウェア工学}

AIとソフトウェア工学の関係は、主に二方向で整理できる。第一に、ソフトウェア工学へのAI応用である。これは、要求の曖昧性解消、欠陥予測、テストケース生成、脆弱性分析、保守性予測、工程評価などにAIを使う考え方である。第二に、AIシステムのためのソフトウェア工学である。AIシステムでは、振る舞いが明示的な規則だけでなく訓練データから得られるため、従来型システムとは異なる開発支援が必要になる。

AIシステム開発では、データ科学者とソフトウェア技術者の協働、大規模で変化するデータセットを前提にした進化、倫理、公平性、説明可能性、データ品質に関する要求が重要になる。機械学習ソフトウェアの設計パターンのように、AIシステム向けの実践を形式化する動きもある。

\begin{importantbox}{注意}
AIを使えばソフトウェア工学上の判断が不要になるわけではない。AIの出力には不確実性があり、十分に構造化されたデータ、評価、監視、失敗時の扱いが必要である。AIシステムほど、要求、品質、セキュリティ、倫理を明示的に扱う必要がある。

\end{importantbox}

\begin{todobox}
\textbf{AI for SE と SE for AI の関係図を入れる。}\par
\textbf{画像生成用プロンプト}\par
白背景、モノクロ、A4本文幅の双方向図を作成する。左に「AI for SE：ソフトウェア工学へのAI応用」を置き、例として「欠陥予測」「テスト生成」「保守性予測」「脆弱性分析」を並べる。右に「SE for AI：AIシステムのためのソフトウェア工学」を置き、例として「データ要求」「倫理・公平性」「モデル監視」「ML設計パターン」を並べる。中央に「相互に支える」と書く。日本語ラベル、講義ノート風。

\end{todobox}

\section{10 トピックと参考文献の対応}

この表は、SWEBOK 16章が示すトピック対参考資料の行列を、学習用に簡略化したものである。詳しく学ぶ場合は、各トピックに対応する文献を読むとよい。

\begin{longtable}{>{\raggedright\arraybackslash}p{0.36\textwidth}>{\raggedright\arraybackslash}p{0.58\textwidth}}
\toprule
トピック & 主な参考資料 \\
\midrule
\endfirsthead
\toprule
トピック & 主な参考資料 \\
\midrule
\endhead
1 システムまたは解決策の基本概念 & Sommerville, Software Engineering。 \\
2 コンピュータアーキテクチャと組織 & Null and Lobur, The Essentials of Computer Organization and Architecture; Sommerville。 \\
3 データ構造とアルゴリズム & Cormen et al., Introduction to Algorithms; Horowitz et al., Computer Algorithms; Null and Lobur。 \\
4 プログラミング基礎と言語 & Brookshear, Computer Science: An Overview; Null and Lobur; McConnell, Code Complete。 \\
5 基本ソフトウェア & Tanenbaum and Bos, Modern Operating Systems; Brookshear。 \\
6 データベース管理 & Connolly and Begg, Database Systems; Hernandez, Database Design for Mere Mortals; Fagin。 \\
7 コンピュータネットワークと通信 & Kurose and Ross, Computer Networking: A Top-Down Approach; Brookshear。 \\
8 人間要因 & Nielsen, Usability Engineering; McConnell; ISO 9241-420。 \\
9 人工知能と機械学習 & Goodfellow et al., Deep Learning; Shafiq et al.; Martínez-Fernández et al.; Washizaki et al.。 \\
\bottomrule
\end{longtable}

\section{11 参考文献}

以下は SWEBOK 16章で参照される文献を、学習用に整理したものである。

\begin{enumerate}[leftmargin=2em]
\item Joint Task Force on Computing Curricula, IEEE Computer Society and Association for Computing Machinery, Software Engineering 2014: Curriculum Guidelines for Undergraduate Degree Programs in Software Engineering, 2014.
\item G. Voland, Engineering by Design, 2nd ed., Prentice Hall, 2003.
\item S. McConnell, Code Complete, 2nd ed., Microsoft Press, 2004.
\item J. G. Brookshear, Computer Science: An Overview, 12th ed., Addison-Wesley, 2017.
\item E. Horowitz et al., Computer Algorithms, 2nd ed., Silicon Press, 2007.
\item I. Sommerville, Software Engineering, 10th ed., Addison-Wesley, 2016.
\item ISO/IEC/IEEE 24765:2017, Systems and Software Engineering — Vocabulary, 2nd ed., 2017.
\item L. Null and J. Lobur, The Essentials of Computer Organization and Architecture, 5th ed., Jones and Bartlett Publishers, 2018.
\item J. Nielsen, Usability Engineering, Morgan Kaufmann, 1994.
\item ISO 9241-420:2011, Ergonomics of Human-System Interaction, ISO, 2011.
\item M. Bishop, Computer Security: Art and Science, 2nd ed., Addison-Wesley, 2018.
\item R. C. Seacord, The CERT C Secure Coding Standard, Addison-Wesley Professional, 2016.
\item R. Fagin, A Normal Form for Relational Databases that is based on Domains and Keys, ACM Transactions on Database Systems, 1981.
\item I. Goodfellow, Y. Bengio, A. Courville, Deep Learning, 2018.
\item S. Shafiq, A. Mashkoor, C. Mayr-Dorn, A. Egyed, A Literature Review of Using Machine Learning in Software Development Life Cycle Stages, IEEE Access, 2021.
\item S. Martínez-Fernández et al., Software Engineering for AI-Based Systems: A Survey, ACM Transactions on Software Engineering and Methodology, 2022.
\item H. Washizaki et al., Software Engineering Design Patterns for Machine Learning Applications, Computer, 2022.
\item T. H. Cormen, C. E. Leiserson, R. L. Rivest, C. Stein, Introduction to Algorithms, 4th ed., 2022.
\item A. S. Tanenbaum and H. Bos, Modern Operating Systems, 4th ed., 2016.
\item IEEE Xplore document 9779481, referenced by SWEBOK 16章 for Harvard architecture.
\item N. Ford, M. Richards, P. Sadalage, Z. Dehghani, Software Architecture: The Hard Parts, O'Reilly, 2021.
\item T. Connolly and C. Begg, Database Systems: A Practical Approach to Design, Implementation and Management, 6th ed., Pearson.
\item M. J. Hernandez, Database Design for Mere Mortals, 4th ed., Addison-Wesley.
\item J. F. Kurose and K. W. Ross, Computer Networking: A Top-Down Approach, 7th ed., Pearson.
\end{enumerate}

\section{12 画像 TODO 一覧}

本文に配置した画像 TODO を再掲する。実際に図を作る場合は、各 TODO のプロンプトをそのまま画像生成に使うと、A4資料に合う図を作りやすい。

\begin{longtable}{>{\raggedright\arraybackslash}p{0.07\textwidth}>{\raggedright\arraybackslash}p{0.25\textwidth}>{\raggedright\arraybackslash}p{0.62\textwidth}}
\toprule
No. & TODO & 画像生成用プロンプト \\
\midrule
\endfirsthead
\toprule
No. & TODO & 画像生成用プロンプト \\
\midrule
\endhead
1 & 16章の全体地図を入れる。 & 白背景、モノクロ、A4本文幅に収まる横長の学習用図解を作成する。中央に「コンピューティング基礎」を置き、周囲に「システム概念」「コンピュータ構成」「データ構造とアルゴリズム」「プログラミング」「基本ソフトウェア」「データベース」「ネットワーク」「人間要因」「人工知能・機械学習」を配置する。矢印で「ソフトウェア技術者の設計判断を支える知識」と示す。ラベルは日本語、細線、講義ノート風、余白を広くする。 \\
2 & システム分解と三つの設計性質を示す図を入れる。 & 白背景、モノクロ、A4本文幅の概念図を作成する。大きな箱を「システム」とし、その中に「サブシステム」「モジュール」「構成要素」を階層的に配置する。横に三つの吹き出しで「モジュール性：適切な単位に分ける」「凝集性：一つの仕事に集中する」「疎結合性：影響範囲を小さくする」を示す。日本語ラベル、細線、講義ノート風。 \\
3 & 代表的なコンピュータ構成の比較図を入れる。 & 白背景、モノクロ、A4本文幅の比較図を作成する。左に「ノイマン型：命令とデータが同じ記憶」、中央に「ハーバード型：命令記憶とデータ記憶を分離」、右に「Flynn分類：SISD/SIMD/MISD/MIMD」を置く。各構成に小さなブロック図を添える。ラベルは日本語、細線、講義ノート風。 \\
4 & CPU、メモリ、入出力、制御装置の関係図を入れる。 & 白背景、モノクロ、A4本文幅のブロック図を作成する。中央にCPUを置き、その中に「算術論理装置」「制御装置」「レジスタ」を配置する。周囲に「主記憶」「入力装置」「出力装置」「補助記憶」を置き、アドレスバス、データバス、制御バスを異なる線種で示す。日本語ラベル、講義ノート風。 \\
5 & データ構造の分類ツリーを入れる。 & 白背景、モノクロ、A4本文幅の階層図を作成する。最上位に「データ構造」を置き、「基本データ型」「複合データ型」「抽象データ型」に分岐する。複合データ型から「線形：配列・リスト・スタック・キュー」「非線形：木・ヒープ・グラフ」に分岐する。日本語ラベル、細線、講義ノート風。 \\
6 & 計算量の増え方を比較するグラフを入れる。 & 白背景、モノクロ、A4本文幅の折れ線グラフを作成する。横軸を「入力サイズ n」、縦軸を「処理量」とし、O(1)、O(log n)、O(n)、O(n log n)、O(n\textasciicircum{}2)、O(2\textasciicircum{}n) の曲線を描く。指数の曲線は急上昇、対数は緩やかにする。日本語ラベル、講義ノート風、凡例付き。 \\
7 & ソート・探索・ハッシュの使い分け図を入れる。 & 白背景、モノクロ、A4本文幅の判断フローを作成する。開始は「データから目的の結果を得たい」。分岐として「順序が必要ならソート」「一件を探すなら探索」「キーで高速アクセスするならハッシュ」を置く。各枝に例として「一覧表示」「ユーザー検索」「辞書・キャッシュ」を添える。日本語ラベル、講義ノート風。 \\
8 & ソースコードが実行されるまでの流れを入れる。 & 白背景、モノクロ、A4本文幅の工程図を作成する。「ソースコード」から「字句解析」「構文解析」「中間表現」「最適化」「コード生成」「リンク」「ロード」「実行」へ矢印でつなぐ。横に「インタプリタでは一部工程が実行時に行われる」と注記する。日本語ラベル、講義ノート風。 \\
9 & サブルーチンとコルーチンの制御の違いを示す図を入れる。 & 白背景、モノクロ、A4本文幅の比較図を作成する。左に「通常のサブルーチン：呼び出し→実行→戻る」を縦の矢印で示す。右に「コルーチン：再開→一時停止→別の処理→再開」を複数の横矢印で示す。日本語ラベル、講義ノート風。 \\
10 & オブジェクト指向の四要素を示す図を入れる。 & 白背景、モノクロ、A4本文幅の4象限図を作成する。象限に「抽象化」「カプセル化」「継承」「多態性」を配置し、それぞれに短い日本語説明を付ける。中央に「オブジェクト＝データ＋操作」を置く。講義ノート風、細線。 \\
11 & OSが管理する資源の全体図を入れる。 & 白背景、モノクロ、A4本文幅の概念図を作成する。中央に「基本ソフトウェア（OS）」を置き、周囲に「CPU」「メモリ」「入出力装置」「ファイル」「ネットワーク」「アプリケーション」を配置する。OSから各資源へ管理矢印を描く。日本語ラベル、講義ノート風。 \\
12 & スキーマとキーの関係図を入れる。 & 白背景、モノクロ、A4本文幅のデータベース図を作成する。表「顧客」と表「注文」を置き、顧客IDを主キー、注文表の顧客IDを外部キーとして矢印でつなぐ。複合キーと代理キーの小さな例を注記する。日本語ラベル、講義ノート風。 \\
13 & 正規化で表が分かれる様子を入れる。 & 白背景、モノクロ、A4本文幅のビフォー・アフター図を作成する。左に「注文表」に顧客名・住所・商品名が重複して入っている例を置く。右に「顧客表」「注文表」「商品表」へ分割した例を置き、主キーと外部キーで関係を示す。日本語ラベル、講義ノート風。 \\
14 & OSI参照モデルと層間通信を示す図を入れる。 & 白背景、モノクロ、A4本文幅の縦長図を作成する。左右に二つのコンピュータを置き、それぞれ7層の積み重ねとして「物理」「データリンク」「ネットワーク」「トランスポート」「セッション」「プレゼンテーション」「アプリケーション」を描く。同じ層同士に点線で「プロトコル」、隣接層に「インタフェース」を示す。日本語ラベル、講義ノート風。 \\
15 & カプセル化とカプセル解除の図を入れる。 & 白背景、モノクロ、A4本文幅の左右対称図を作成する。左側を送信側として、上位データに各層のヘッダが順に付く様子を下向きに描く。右側を受信側として、各層でヘッダが外れる様子を上向きに描く。中央に「PDU」を置く。日本語ラベル、講義ノート風。 \\
16 & 利用者要因と開発者要因を対比する図を入れる。 & 白背景、モノクロ、A4本文幅の2列図を作成する。左列を「利用者要因」とし、「直感的操作」「明確なメッセージ」「高速応答」「取り消し」「安全性」を置く。右列を「開発者要因」とし、「可読性」「文書化」「命名規則」「モジュール化」「保守性」を置く。中央に「人間の限界を前提に設計する」と書く。日本語ラベル、講義ノート風。 \\
17 & 学習方法と代表モデルの対応図を入れる。 & 白背景、モノクロ、A4本文幅のマップ図を作成する。左に「教師あり」「教師なし」「半教師あり」「強化学習」を置き、右に「回帰」「分類」「クラスタリング」「意思決定」「制御」を置く。代表モデルとして線形回帰、決定木、ニューラルネットワーク、ランダムフォレストを小さなタグで配置する。日本語ラベル、講義ノート風。 \\
18 & AI for SE と SE for AI の関係図を入れる。 & 白背景、モノクロ、A4本文幅の双方向図を作成する。左に「AI for SE：ソフトウェア工学へのAI応用」を置き、例として「欠陥予測」「テスト生成」「保守性予測」「脆弱性分析」を並べる。右に「SE for AI：AIシステムのためのソフトウェア工学」を置き、例として「データ要求」「倫理・公平性」「モデル監視」「ML設計パターン」を並べる。中央に「相互に支える」と書く。日本語ラベル、講義ノート風。 \\
\bottomrule
\end{longtable}

\section{13 章末確認問題}

\begin{enumerate}[leftmargin=2em]
\item ソフトウェア技術者と単なるプログラマの役割の違いを説明せよ。
\item モジュール性、凝集性、疎結合性の違いを、身近な例で説明せよ。
\item コンピュータアーキテクチャとコンピュータ組織の違いを説明せよ。
\item ノイマン型とハーバード型の違いを説明せよ。
\item RISC と CISC の違いを、命令の複雑さと実行時間の観点から説明せよ。
\item 線形データ構造と非線形データ構造の違いを説明し、それぞれ例を二つ挙げよ。
\item 大O記法が実行時間そのものではなく増え方を表す理由を説明せよ。
\item サブルーチン、再帰、コルーチンの違いを説明せよ。
\item プロセスとスレッド、プロセス間通信、デッドロックの関係を説明せよ。
\item ACID と BASE の違いを、どのようなデータベースに向くかを含めて説明せよ。
\item OSI参照モデルの七層を順に挙げ、カプセル化とは何かを説明せよ。
\item AI for SE と SE for AI の違いを説明し、それぞれ例を二つ挙げよ。
\end{enumerate}

\begin{importantbox}{解答の方向性}
1は、実装だけでなく要求、設計、性能、テスト、保守まで扱う点を述べる。2は、分割、責務集中、依存低減を区別する。3は、外から見た構成と内部実現の違いを述べる。4から12は、本文の定義、表、注意欄の語を使って説明できる。

\end{importantbox}

\end{document}
